% Source: https://artofproblemsolving.com/community/c3h1317067_temperatures
% Stable identity: aops-topic-1317067-q01
$F=\dfrac{9}{5}C+32$. For $C=\overline{ab}$ we need $F$ to round to $\overline{ba}$. Since $F=1.8C+32$ grows faster than $C$, a two-digit $C$ with a two-digit reversed $F$ requires $F>C$, so $b>a$; also $F$ has two digits only for $C\le37$ (as $C=38$ gives $F=100.4$). Checking $C=10,\dots,37$:
\begin{itemize}
\item $C=16$: $F=60.8\approx61$, reversed digits. \checkmark
\item $C=28$: $F=82.4\approx82$, reversed digits. \checkmark
\end{itemize}
and no other value works. So the temperatures are $\boxed{16^\circ\text{C}\ (61^\circ\text{F})\text{ and }28^\circ\text{C}\ (82^\circ\text{F})}$. (No two-digit temperature works exactly: $F$ is an integer only when $5\mid C$, and none of $10,15,\dots,35$ works.)
